
Projected on the numerical distribution function grid, and taking into
account of the internal and external boundaries, the $3-D$ electron drift
kinetic equation may be expressed as 
\begin{eqnarray}
&&\sum_{i^{\prime }=i-1}^{i^{\prime }=i+1}\sum_{j^{\prime \prime }=j^{\prime
}-1}^{j^{\prime \prime }=j^{\prime }+1}\widehat{\overline{M}}%
_{p,l+1/2,i^{\prime }+1/2,j^{\prime }+1/2}^{\left( 0\right)
}g_{l+1/2,i^{\prime }+1/2,j^{\prime \prime }+1/2}^{\left( 0\right) (k)}= 
\nonumber \\
&&\widehat{\widetilde{W}}_{p,l+1/2,i+1/2,j^{\prime }+1/2}^{\left( 0\right)
}-p_{i+1/2}^{2}C^{\left( 0\right) }\left( f_{M,l+1/2,i+1/2,j^{\prime }+1/2},%
\frac{3}{2}\xi _{0,j+1/2}g_{l+1/2,i+1/2,j^{\prime }+1/2}^{\left( 0\right)
(m=1)(k)}\right)  \label{DKeq1}
\end{eqnarray}
where 
\begin{eqnarray}
\widehat{\widetilde{W}}_{p,l+1/2,i+1/2,j+1/2}^{\left( 0\right) }
&=&-\sum_{i^{\prime }=i-1}^{i^{\prime }=i+1}\sum_{j^{\prime \prime
}=j-1}^{j^{\prime \prime }=j+1}\widetilde{\overline{M}}_{p,l+1/2,i^{\prime
}+1/2,j^{\prime \prime }+1/2}^{\left( 0\right) }\widetilde{f}%
_{l+1/2,i^{\prime }+1/2,j^{\prime \prime }+1/2}^{\left( 0\right) }  \nonumber
\\
&&-\sum_{i^{\prime }=i-1}^{i^{\prime }=i+1}\sum_{j^{\prime \prime
}=j-1}^{j^{\prime \prime }=j+1}\widetilde{\overline{H}}_{\psi
,l+1/2,i^{\prime }+1/2,j^{\prime \prime }+1/2}^{\left( 0\right) }\widetilde{f%
}_{l+1/2,i^{\prime }+1/2,j^{\prime \prime }+1/2}^{\left( 0\right) } 
\nonumber \\
&&-p_{i+1/2}^{2}\widetilde{C}^{\left( 0\right) }\left(
f_{M,l+1/2,i+1/2,j+1/2},\frac{3}{2}\xi _{0,j+1/2}\widetilde{f}%
_{l+1/2,i+1/2,j+1/2}^{\left( 0\right) (m=1)}\right)
\end{eqnarray}

As mentioned in Sec.\ref{dkeboundaries}, $j\rightarrow \left\{ 0,n_{\xi
_{0}}-1\right\} ,$ while $j^{\prime }\rightarrow \left\{ 0,n_{\xi
_{0}}-m_{\xi _{0T,l+1/2}}-1\right\} $ corresponds to the reduced pitch-angle
grid which depends of the radial position $\psi _{l+1/2}$, since all the
trapped region has been removed. In Eq.\ref{DKeq1}, time evolution is not
considered, since function $g^{\left( 0\right) }$ is only determined for
steady-state value of $\widetilde{f}^{\left( 0\right) }.$ Therefore,
iterations to reach the solution only result from the non-linear electron
self-collision operator $C^{\left( 0\right) }$, that ensures momentum
conservation.

In compact matrix form, the differential equation in the reduced momentum
phase space has the following symbolic form 
\begin{equation}
\widehat{\Bbb{G}}Y^{(k+1)}=\widehat{\Bbb{G}}Y^{(k)}+\widehat{C}\left(
X_{M},Y^{(k)}\right) +\widehat{\widetilde{W}}  \label{Matrixeqg}
\end{equation}
where $\widehat{\Bbb{G}}$ corresponds to the flux divergence, $\widehat{C}$
is the first order Legendre correction and $\widehat{\widetilde{W}}$ to the
contribution of the radial gradient of the distribution function $%
f_{0}^{\left( 0\right) }$ . Here $Y^{(k)}$ is a vector whose components are
the discrete values of the distribution function $g^{\left( 0\right) \left(
k\right) }$ at iteration step $\left( k\right) ,$ organized as follows 
\begin{equation}
Y^{(k)}\rightarrow \left\{ 
\begin{array}{l}
l=0\rightarrow \left\{ 
\begin{array}{l}
i=0\rightarrow \left\{ 
\begin{array}{l}
j^{\prime }=0 \\ 
j^{\prime }=\ldots \\ 
j^{\prime }=n_{\xi _{0}}-m_{\xi _{0T,l+1/2}}-1
\end{array}
\right. \\ 
i=\ldots \\ 
i=n_{p}-1
\end{array}
\right. \\ 
l=\ldots \\ 
l=n_{\psi }-1
\end{array}
\right.
\end{equation}

By construction, $\widehat{\Bbb{G}}$ is simply a block diagonal matrix. Each
block, which describes the momentum dynamics at location $\psi _{l+1/2},$ is
a square matrice of $9$ diagonals whose size $\left( n_{\xi _{0}}-m_{\xi
_{0T,l+1/2}}\right) n_{p}\times $ $\left( n_{\xi _{0}}-m_{\xi
_{0T,l+1/2}}\right) n_{p}$ progressively decreases from $\psi _{1/2}$ to $%
\psi _{n_{\xi _{0}}-1/2}$ because of the trapped/passing boundary enlarges
from the center to the plasma edge, as shown in Fig. \ref{dkematrix}$.$ The
main diagonal of $\widehat{\Bbb{G}}$ is also dominant because collisions
predominate over other physical processes at each plasma radius.

%TCIMACRO{
%\TeXButton{Fig_dkematrix}{\fig{figures/dkematrix.eps}{Qualitative shape of the matrix for the drift kinetic equation}{dkematrix}{10cm}}}
%BeginExpansion
\fig{figures/dkematrix.eps}{Qualitative shape of the matrix for the drift kinetic equation}{dkematrix}{10cm}%
%EndExpansion

The matrix $\widehat{\Bbb{G}}$ has a global banded structure with an
extremely high sparcity, a property that is extensively used for the storage
in the computer memory, like for the Fokker-Planck equation. Indeed, the
size of $\widehat{\Bbb{G}}$ is 
\begin{equation}
size\left( \widehat{\Bbb{G}}\right) =n_{p}^{2}\left[ \sum_{l=0}^{l=n_{\psi
}-1}\left( n_{\xi _{0}}-m_{\xi _{0T,l+1/2}}\right) \right] ^{2}\leq n_{\psi
}^{2}n_{p}^{2}n_{\xi _{0}}^{2}
\end{equation}

For tokamak with very large aspect ratio, $size\left( \widehat{\Bbb{G}}%
\right) \approx n_{\psi }^{2}n_{p}^{2}n_{\xi _{0}}^{2}$ since trapped
particle contribution is fairly negligible. In this limit, $size\left( 
\widehat{\Bbb{G}}\right) \simeq size\left( \widehat{\Bbb{B}}\right) ,$ where 
$\widehat{\Bbb{B}}$ is the flux matrix of the Fokker-Planck equation. For
the values of $n_{\psi }$, $n_{p}$ and $n_{\xi _{0}}$ taken in Sec. \ref
{fpmatrixsection}, the required memory is approximately $57.6MBytes$, a
lower level than for the Fokker-Planck equation.
